\section{Operational readiness}
\label{sec:operations}

\begin{question}
  \textbf{Is the system ready for production training?} How long until training starts, how often do runs
  fail, how much do repeated runs vary, and is throughput stable over an hour or more?
\end{question}

\paragraph{Startup.} \Cref{fig:startup} separates process start to an initialized model (imports, process
groups, model construction and sharding) from the first step (compilation, or loading the compile cache).

\begin{figure}[h]
  \centering\fig{startup}
  \caption{Startup phases by number of nodes. \syntheticnote{startup}}
  \label{fig:startup}
\end{figure}

\placeholder{outliers are explained: 297\,s init with the virtual environment on Lustre instead of node-local
disk; first steps of 85--117\,s with a cold compile cache (new shapes).}

\paragraph{Stability.} \Cref{fig:stability} follows a long run step by step, with the GPU clocks from DCGM on
the same time axis, so that a throughput drop can be traced to a GPU.

\begin{figure}[h]
  \centering\fig{stability}
  \caption{\syntheticnote{stability}Throughput per step over one hour on every healthy node (top) and SM clocks of
    all GPUs (bottom): the gray band is the range of the other GPUs, the colored line the GPU with the
    lowest clock.}
  \label{fig:stability}
\end{figure}

\paragraph{Stragglers.} With FSDP, all ranks wait for the slowest one, so per-rank step times are identical
(measured: 24 ranks within 0.2\,ms) and cannot identify a slow GPU. The slow GPU is instead the one that
\emph{waits least}: its tensor cores are busy for a larger share of each step (\cref{fig:stragglers}).

\begin{figure}[h]
  \centering\fig{stragglers}
  \caption{\syntheticnote{stragglers}Tensor-pipe activity per GPU (DCGM), the highest 10-minute median of each GPU,
    by leaf. The outlier is named.}
  \label{fig:stragglers}
\end{figure}

\placeholder{run-to-run variance (repetitions in \cref{fig:placement-scaling,fig:placement-at-scale}; so far $\le$0.1\,\% on Rorqual within an
allocation), failures and their causes (\cref{tab:ledger}), excluded nodes.}
